% ============================== 4. ИМПЛЕМЕНТАЦИЯ И РЕЗУЛТАТИ ==============================
\chapter{Имплементация и експериментални резултати}\label{ch:results}

Главата представя резултатите в реда на експерименталния конвейер: данни и обработка, фаза на машинно обучение и проверка за преобучение, оперативна фаза, статистическа надеждност, ефект на редукцията и работа при естествено съотношение. Всички числа, таблици и фигури са генерирани автоматично от последното изпълнение на \texttt{run\_all.py}~\cite{repo}.

% ---------------------------------------------------------------------------
\section{Данни и предварителна обработка}\label{sec:r-data}

От изтеглените \val{rows-read} записа \val{nan-rows} съдържат липсващи или безкрайни стойности, а \val{dup-rows} (\val{dup-share}) са пълни дубликати. След почистването остават \val{rows-clean} уникални записа (\cref{tab:class-counts}). Почистването на ниво задача (\cref{eq:conflict}) премахва още \val{task-conf-Bot} записа (\val{task-confvec-Bot} вектора) с противоречиви етикети в задачата Bot, \val{task-conf-DDoS} в задачата DDoS и \val{task-dup-BruteForce} дубликата в задачата BruteForce. От всеки клас се изтеглят по \val{per-class-DDoS} записа за DDoS и Bot и \val{per-class-BruteForce} за BruteForce, т.е. всички уникални записи на BruteForce. Така обучаващите множества съдържат \val{train-rows-DDoS}, \val{train-rows-BruteForce} и \val{train-rows-Bot} записа, а тестовите \val{test-rows-DDoS}, \val{test-rows-BruteForce} и \val{test-rows-Bot} (\cref{fig:classes}).

\begin{figure}[H]
\centering
\includegraphics[width=\textwidth]{01_class_distribution.png}
\caption{Разпределение на класовете в изтеглените данни и след балансирането (логаритмична скала)}\label{fig:classes}
\vspace{2pt}
\interpret{Сивите стълбове са нормалният трафик, а червените атаката. В оригиналните файлове нормалният трафик е 2,6 (DDoS) до 6,3 пъти (Bot) повече от атаката. Вдясно е балансираната извадка 1:1, върху която се обучават и тестват моделите. Оригиналното съотношение се използва отново в експеримент E2.}
\end{figure}

\Cref{fig:ks} подрежда признаците по KS статистиката от \cref{eq:ks} (анализ A1). При BruteForce един признак, \texttt{\val{ks-top-BruteForce}}, сам разделя класовете ($D=\val{ks-topD-BruteForce}$), а средното $D$ на петте най-добри признака е \val{ks-top5-BruteForce}. При DDoS най-разделим е \texttt{\val{ks-top-DDoS}} ($D=\val{ks-topD-DDoS}$), а при Bot \texttt{\val{ks-top-Bot}} ($D=\val{ks-topD-Bot}$). Признаци с $D>0{,}5$ има \val{ks-n05-DDoS} за DDoS, \val{ks-n05-BruteForce} за BruteForce и \val{ks-n05-Bot} за Bot.

\begin{figure}[H]
\centering
\includegraphics[width=\textwidth]{12_ks_separability.png}
\caption{Едномерна разделимост на признаците: KS статистика $D$ между атаката и нормалния трафик}\label{fig:ks}
\vspace{2pt}
\interpret{Дължината на стълба е максималната разлика между кумулативните разпределения на признака при атака и при нормален трафик. Стойност близо до 1 означава, че само този признак почти разделя класовете. При BruteForce много признаци имат $D>0{,}9$: опитите за вход са кратки потоци с еднаква структура. При DDoS и Bot най-добрите признаци са около 0,75, т.е. всеки сам по себе си не стига, но заедно разделят класовете.}
\end{figure}

Редукцията свива вектора на признаците в два етапа, както е обобщено в \cref{tab:reduction}. Филтърът премахва \val{lowvar-DDoS} константни признака (\texttt{\val{lowvar-names}}), а подборът запазва \val{n-selected} от останалите \val{n-filtered}. Избраните признаци носят \val{sel-share-DDoS} (DDoS), \val{sel-share-BruteForce} (BruteForce) и \val{sel-share-Bot} (Bot) от общата важност в помощния модел. Най-важни са \texttt{\val{sel-top3-DDoS}} за DDoS, \texttt{\val{sel-top3-BruteForce}} за BruteForce и \texttt{\val{sel-top3-Bot}} за Bot.

\begin{table}[H]
\centering\small
\caption{Редукция на размерността по задача: филтър и подбор на признаци}\label{tab:reduction}
\fitbox{%
\begin{tabular}{@{}lccccc>{\raggedright\arraybackslash}p{6.2cm}@{}}
\toprule
\headrow \hd{Задача} & \hd{Вход} & \hd{Нулева\\дисперсия} & \hd{След\\филтър} & \hd{Избрани} & \hd{Дял на\\важността} & \hd{Избрани признаци}\\
\midrule
DDoS & 77 & 10 & 67 & 20 & 99,9\% & \scriptsize\texttt{Protocol, Flow Duration, Tot Fwd Pkts, Tot Bwd Pkts, TotLen Bwd Pkts, Bwd Pkt Len Mean, Bwd Pkt Len Std, Flow Pkts/s, Flow IAT Max, Flow IAT Min, Fwd IAT Tot, Fwd IAT Max, Fwd IAT Min, Fwd Header Len, Fwd Pkts/s, Pkt Len Mean, Pkt Len Std, Init Fwd Win Byts, Init Bwd Win Byts, Fwd Act Data Pkts}\\
\rowcolor{rowalt}BruteForce & 77 & 10 & 67 & 20 & 100,0\% & \scriptsize\texttt{Flow Duration, Tot Fwd Pkts, Tot Bwd Pkts, TotLen Bwd Pkts, Bwd Pkt Len Max, Bwd Pkt Len Mean, Flow Pkts/s, Flow IAT Mean, Flow IAT Max, Flow IAT Min, Fwd IAT Min, Bwd IAT Mean, Bwd IAT Min, Fwd Header Len, Fwd Pkts/s, Pkt Len Std, Init Fwd Win Byts, Init Bwd Win Byts, Fwd Act Data Pkts, Fwd Seg Size Min}\\
Bot & 77 & 10 & 67 & 20 & 99,5\% & \scriptsize\texttt{Protocol, Tot Fwd Pkts, Fwd Pkt Len Max, Fwd Pkt Len Mean, Bwd Pkt Len Max, Bwd Pkt Len Mean, Flow Byts/s, Flow IAT Std, Flow IAT Max, Flow IAT Min, Fwd IAT Std, Fwd IAT Min, Fwd Pkts/s, Pkt Len Max, Pkt Len Mean, RST Flag Cnt, Pkt Size Avg, Init Fwd Win Byts, Active Mean, Idle Min}\\
\bottomrule
\end{tabular}}%

\vspace{2pt}
\interprettab{За всяка задача е показан броят признаци на входа, броят премахнати като константни, броят след филтъра, броят избрани признаци и делът на общата важност, който те носят. Векторът се свива от \val{n-features} до \val{n-selected} измерения (с 74\%), като почти цялата информация се запазва. Подборът се извършва отделно за всяка задача, защото различните атаки се разпознават по различни признаци.}
\end{table}

Двумерните проекции PCA и t-SNE~\cite{vandermaaten2008} (анализ A2) показват, че и при трите атаки атакуващите записи образуват компактни групи, отделени от нормалния трафик. За контролния метод PCA 95\% от дисперсията се запазват с \val{pca-DDoS} компонента за DDoS, \val{pca-BruteForce} за BruteForce и \val{pca-Bot} за Bot (\cref{eq:kstar}).

% ---------------------------------------------------------------------------
\section{Фаза на машинно обучение и проверка за преобучение}\label{sec:r-ml}

\Cref{tab:cv} показва избраните хиперпараметри. Правилото „една стандартна грешка“ (\cref{eq:one-se}) избира по-проста конфигурация от тази с най-висока оценка в няколко случая (например SVM за Bot с $C=1$ вместо $C=10$), без статистически значима загуба. Разсейването на F1 между петте фолда (експеримент E3) е в рамките на стотни от процента при DDoS и BruteForce и около една десета от процента при Bot.

\begin{table}[H]
\centering\small
\caption{Резултат от GridSearchCV: избрани хиперпараметри и CV F1}\label{tab:cv}
\fitbox{%
\begin{tabular}{@{}llp{4.4cm}rrrrr@{}}
\toprule
\headrow \hd{Модел} & \hd{Атака} & \hd{Избрани хиперпараметри} & \hd{CV F1\\избрана} & \hd{CV F1\\най-висока} & \hd{Ст. откл.} & \hd{Мин.--макс. по фолд} & \hd{Време, s}\\
\midrule
Random Forest & DDoS & \scriptsize\texttt{max\_depth=25, max\_features=0.5, min\_samples\_leaf=2, n\_estimators=200} & 99,99\% & 99,99\% & 0,01\% & 99,98\% -- 100,00\% & 241\\
\rowcolor{rowalt}SVM (RBF) & DDoS & \scriptsize\texttt{C=10, gamma=0.1} & 99,96\% & 99,97\% & 0,03\% & 99,90\% -- 100,00\% & 12\\
XGBoost & DDoS & \scriptsize\texttt{learning\_rate=0.03, max\_depth=6, min\_child\_weight=1, n\_estimators=300} & 100,00\% & 100,00\% & 0,00\% & 100,00\% -- 100,00\% & 61\\
\rowcolor{rowalt}Random Forest & BruteForce & \scriptsize\texttt{max\_depth=25, max\_features=sqrt, min\_samples\_leaf=2, n\_estimators=200} & 99,99\% & 99,99\% & 0,01\% & 99,97\% -- 100,00\% & 189\\
SVM (RBF) & BruteForce & \scriptsize\texttt{C=10, gamma=scale} & 99,96\% & 99,97\% & 0,03\% & 99,92\% -- 100,00\% & 10\\
\rowcolor{rowalt}XGBoost & BruteForce & \scriptsize\texttt{learning\_rate=0.03, max\_depth=6, min\_child\_weight=5, n\_estimators=300} & 99,99\% & 99,99\% & 0,01\% & 99,98\% -- 100,00\% & 53\\
Random Forest & Bot & \scriptsize\texttt{max\_depth=25, max\_features=sqrt, min\_samples\_leaf=2, n\_estimators=200} & 99,93\% & 99,94\% & 0,03\% & 99,90\% -- 99,98\% & 226\\
\rowcolor{rowalt}SVM (RBF) & Bot & \scriptsize\texttt{C=1, gamma=scale} & 99,47\% & 99,49\% & 0,10\% & 99,35\% -- 99,60\% & 15\\
XGBoost & Bot & \scriptsize\texttt{learning\_rate=0.1, max\_depth=6, min\_child\_weight=5, n\_estimators=800} & 99,96\% & 99,97\% & 0,02\% & 99,94\% -- 99,99\% & 70\\
\bottomrule
\end{tabular}}%
\par\smallskip{\footnotesize\color{muted}Избраната конфигурация е най-простата в рамките на една стандартна грешка от най-високата CV оценка (правило 1-SE).}

\vspace{2pt}
\interprettab{За всеки модел е посочена избраната конфигурация, нейното средно F1 по петте фолда, най-високото F1 в решетката, стандартното отклонение, най-ниската и най-високата стойност по фолд и времето за търсене. Всички модели постигат CV F1 над 99,4\%. Най-ниско е F1 на SVM при Bot (\val{SVM-Bot-cv}).}
\end{table}

\Cref{tab:overfit} проверява явно за преобучение. Преобученият модел има много по-висока точност върху обучаващите данни, отколкото върху новите. Тук разликата е най-много \val{gap-max} процентни пункта, а F1 от кръстосаната валидация съвпада с F1 върху тестовите данни. Следователно моделите обобщават, а не запомнят обучаващите записи.

\begin{table}[H]
\centering\small
\caption{Проверка за преобучение: обучение, кръстосана валидация и тест}\label{tab:overfit}
\fitbox{%
\begin{tabular}{@{}llrrrrr@{}}
\toprule
\headrow \hd{Модел} & \hd{Атака} & \hd{Accuracy\\обучение} & \hd{F1\\5-CV} & \hd{Accuracy\\тест} & \hd{F1\\тест} & \hd{Разлика обуч.$-$тест,\\пр. п.}\\
\midrule
Random Forest & DDoS & 100,000\% & 99,994\% & 100,000\% & 100,000\% & $+$0,000\\
\rowcolor{rowalt}SVM (RBF) & DDoS & 99,954\% & 99,955\% & 99,950\% & 99,950\% & $+$0,004\\
XGBoost & DDoS & 100,000\% & 100,000\% & 100,000\% & 100,000\% & $+$0,000\\
\rowcolor{rowalt}Random Forest & BruteForce & 99,998\% & 99,992\% & 99,998\% & 99,998\% & $+$0,000\\
SVM (RBF) & BruteForce & 99,972\% & 99,965\% & 99,986\% & 99,986\% & $-$0,014\\
\rowcolor{rowalt}XGBoost & BruteForce & 99,997\% & 99,990\% & 99,998\% & 99,998\% & $-$0,001\\
Random Forest & Bot & 99,983\% & 99,932\% & 99,967\% & 99,967\% & $+$0,016\\
\rowcolor{rowalt}SVM (RBF) & Bot & 99,586\% & 99,473\% & 99,598\% & 99,597\% & $-$0,012\\
XGBoost & Bot & 99,991\% & 99,960\% & 99,982\% & 99,982\% & $+$0,010\\
\bottomrule
\end{tabular}}%

\vspace{2pt}
\interprettab{Колоните сравняват Accuracy върху обучаващите данни, F1 от 5-кратната кръстосана валидация (данни, невиждани при обучението на съответния фолд) и Accuracy и F1 върху тестовите данни от оперативната фаза. Отрицателна разлика означава, че моделът е малко по-точен върху тестовите данни, отколкото върху обучаващите, което е в рамките на случайността.}
\end{table}

Кривите на обучение (анализ A3, \texttt{results/learning\_curves.csv}) показват, че още с \val{lc-RF-DDoS-nfirst} записа всички модели надхвърлят 99,3\%. С повече данни F1 при Bot расте до \val{lc-RF-Bot-last} (RF) и \val{lc-XGB-Bot-last} (XGBoost), а SVM расте по-бавно (\val{lc-SVM-Bot-first} до \val{lc-SVM-Bot-last}) и времето му за обучение расте по-бързо от линейно.

% ---------------------------------------------------------------------------
\section{Оперативна фаза: основни резултати}\label{sec:r-op}

\Cref{tab:main} съдържа всички величини от матрицата на объркване и метриките (\ref{eq:accuracy})--(\ref{eq:spec}) на деветте модела. Те са изчислени от записаните модели върху тестови записи, които не са участвали в обучението.

\begin{table}[H]
\centering\small
\caption{Основни резултати в оперативната фаза}\label{tab:main}
\fitbox{%
\begin{tabular}{@{}llrrrrrrrrr@{}}
\toprule
\headrow \hd{Модел} & \hd{Атака} & \hd{TP} & \hd{TN} & \hd{FP} & \hd{FN} & \hd{Accuracy} & \hd{Precision} & \hd{Recall\\(Sensitivity)} & \hd{F1} & \hd{Specificity}\\
\midrule
Random Forest & PortScan & 504 & 504 & 1 & 1 & 99,80\% & 99,80\% & 99,80\% & 99,80\% & 99,80\%\\
\rowcolor{rowalt}SVM (RBF) & PortScan & 501 & 503 & 2 & 4 & 99,41\% & 99,60\% & 99,21\% & 99,40\% & 99,60\%\\
XGBoost & PortScan & 504 & 503 & 2 & 1 & 99,70\% & 99,60\% & 99,80\% & 99,70\% & 99,60\%\\
\rowcolor{rowalt}Random Forest & DDoS & 30\,000 & 30\,000 & 0 & 0 & 100,00\% & 100,00\% & 100,00\% & 100,00\% & 100,00\%\\
SVM (RBF) & DDoS & 29\,988 & 29\,982 & 18 & 12 & 99,95\% & 99,94\% & 99,96\% & 99,95\% & 99,94\%\\
\rowcolor{rowalt}XGBoost & DDoS & 30\,000 & 30\,000 & 0 & 0 & 100,00\% & 100,00\% & 100,00\% & 100,00\% & 100,00\%\\
Random Forest & BruteForce & 28\,227 & 28\,229 & 0 & 1 & 100,00\% & 100,00\% & 100,00\% & 100,00\% & 100,00\%\\
\rowcolor{rowalt}SVM (RBF) & BruteForce & 28\,221 & 28\,228 & 1 & 7 & 99,99\% & 100,00\% & 99,98\% & 99,99\% & 100,00\%\\
XGBoost & BruteForce & 28\,227 & 28\,229 & 0 & 1 & 100,00\% & 100,00\% & 100,00\% & 100,00\% & 100,00\%\\
\bottomrule
\end{tabular}}%

\vspace{2pt}
\interprettab{TP и TN са верните решения, FP са фалшивите аларми, а FN са пропуснатите атаки. Всички модели имат Accuracy над 99,5\%. Random Forest и XGBoost не правят нито една грешка при DDoS и по една грешка при BruteForce. Най-много грешки има SVM при Bot: \val{SVM-Bot-FN} пропуснати атаки и \val{SVM-Bot-FP} фалшиви аларми от \val{test-rows-Bot} записа.}
\end{table}

\Cref{tab:per100} превежда метриките в броя случаи от 100, както препоръчва учебният материал~\cite{borovska2022}. Това тълкуване е по-разбираемо за оператора на IDS от процентите.

\begin{table}[H]
\centering\scriptsize
\caption{Тълкуване на метриките „на 100 случая“}\label{tab:per100}
\fitbox{%
\begin{tabular}{@{}llrrrrrrrrr@{}}
\toprule
\headrow \hd{Метрика} & \hd{Какво означава} & \hd{Random Forest PortScan} & \hd{SVM (RBF) PortScan} & \hd{XGBoost PortScan} & \hd{Random Forest DDoS} & \hd{SVM (RBF) DDoS} & \hd{XGBoost DDoS} & \hd{Random Forest BruteForce} & \hd{SVM (RBF) BruteForce} & \hd{XGBoost BruteForce}\\
\midrule
Accuracy & правилни решения от 100 тестови записа & 100 от 100 & 99 от 100 & 100 от 100 & 100 от 100 & 100 от 100 & 100 от 100 & 100 от 100 & 100 от 100 & 100 от 100\\
\rowcolor{rowalt}Precision & истински атаки от 100 вдигнати аларми & 100 от 100 & 100 от 100 & 100 от 100 & 100 от 100 & 100 от 100 & 100 от 100 & 100 от 100 & 100 от 100 & 100 от 100\\
Recall & открити атаки от 100 реални атаки & 100 от 100 & 99 от 100 & 100 от 100 & 100 от 100 & 100 от 100 & 100 от 100 & 100 от 100 & 100 от 100 & 100 от 100\\
\rowcolor{rowalt}Specificity & нормални записи без аларма от 100 нормални & 100 от 100 & 100 от 100 & 100 от 100 & 100 от 100 & 100 от 100 & 100 от 100 & 100 от 100 & 100 от 100 & 100 от 100\\
\bottomrule
\end{tabular}}%

\vspace{2pt}
\interprettab{При закръгляне до цял случай от 100 всички модели отговарят правилно в практически всичките 100 случая. Единственото изключение е SVM при Bot: от 100 атаки на Bot \val{SVM-Bot-Recall-100} се откриват.}
\end{table}

\begin{figure}[H]
\centering
\includegraphics[width=0.86\textwidth]{04_confusion_matrices.png}
\caption{Матрици на объркване на деветте модела в оперативната фаза}\label{fig:cm}
\vspace{2pt}
\interpret{Редовете са реалният клас, а колоните предсказаният. Всяка клетка съдържа броя записи и процента от реалния клас. Диагоналът е практически 100\% навсякъде. Грешките на RF и XGBoost се броят на пръсти. SVM при Bot пропуска \val{SVM-Bot-FN} записа (\val{SVM-Bot-FNR} от атаките).}
\end{figure}

ROC кривите на RF и XGBoost са на горния ляв ръб дори при Bot (AUC \val{RF-Bot-AUC} и \val{XGB-Bot-AUC}), а SVM при Bot достига пълен Recall малко по-късно (AUC \val{SVM-Bot-AUC}).

Скоростта в експлоатация се различава съществено. При Bot XGBoost класифицира \val{XGB-Bot-throughput} записа в секунда, Random Forest \val{RF-Bot-throughput}, а SVM \val{SVM-Bot-throughput}. Класификацията със SVM изисква изчисляване на ядрото (\cref{eq:rbf}) спрямо всеки опорен вектор~\cite{chang2011}, а XGBoost използва малко на брой плитки дървета.

% ---------------------------------------------------------------------------
\section{Статистическа надеждност (A9)}\label{sec:r-stat}

\begin{table}[H]
\centering\small
\caption{Метрики с 95\% bootstrap доверителни интервали}\label{tab:ci}
\fitbox{%
\begin{tabular}{@{}llcccc@{}}
\toprule
\headrow \hd{Модел} & \hd{Атака} & \hd{Precision} & \hd{Accuracy} & \hd{Sensitivity} & \hd{F1}\\
\midrule
Random Forest & PortScan & \makecell{99,80\%\\[-1pt]{\scriptsize\color{muted}[99,37\%; 100,00\%]}} & \makecell{99,80\%\\[-1pt]{\scriptsize\color{muted}[99,50\%; 100,00\%]}} & \makecell{99,80\%\\[-1pt]{\scriptsize\color{muted}[99,39\%; 100,00\%]}} & \makecell{99,80\%\\[-1pt]{\scriptsize\color{muted}[99,50\%; 100,00\%]}}\\
\rowcolor{rowalt}SVM (RBF) & PortScan & \makecell{99,60\%\\[-1pt]{\scriptsize\color{muted}[98,99\%; 100,00\%]}} & \makecell{99,41\%\\[-1pt]{\scriptsize\color{muted}[98,91\%; 99,80\%]}} & \makecell{99,21\%\\[-1pt]{\scriptsize\color{muted}[98,37\%; 99,81\%]}} & \makecell{99,40\%\\[-1pt]{\scriptsize\color{muted}[98,90\%; 99,81\%]}}\\
XGBoost & PortScan & \makecell{99,60\%\\[-1pt]{\scriptsize\color{muted}[98,99\%; 100,00\%]}} & \makecell{99,70\%\\[-1pt]{\scriptsize\color{muted}[99,41\%; 100,00\%]}} & \makecell{99,80\%\\[-1pt]{\scriptsize\color{muted}[99,39\%; 100,00\%]}} & \makecell{99,70\%\\[-1pt]{\scriptsize\color{muted}[99,37\%; 100,00\%]}}\\
\rowcolor{rowalt}Random Forest & DDoS & \makecell{100,00\%\\[-1pt]{\scriptsize\color{muted}[100,00\%; 100,00\%]}} & \makecell{100,00\%\\[-1pt]{\scriptsize\color{muted}[100,00\%; 100,00\%]}} & \makecell{100,00\%\\[-1pt]{\scriptsize\color{muted}[100,00\%; 100,00\%]}} & \makecell{100,00\%\\[-1pt]{\scriptsize\color{muted}[100,00\%; 100,00\%]}}\\
SVM (RBF) & DDoS & \makecell{99,94\%\\[-1pt]{\scriptsize\color{muted}[99,91\%; 99,96\%]}} & \makecell{99,95\%\\[-1pt]{\scriptsize\color{muted}[99,93\%; 99,97\%]}} & \makecell{99,96\%\\[-1pt]{\scriptsize\color{muted}[99,94\%; 99,98\%]}} & \makecell{99,95\%\\[-1pt]{\scriptsize\color{muted}[99,93\%; 99,97\%]}}\\
\rowcolor{rowalt}XGBoost & DDoS & \makecell{100,00\%\\[-1pt]{\scriptsize\color{muted}[100,00\%; 100,00\%]}} & \makecell{100,00\%\\[-1pt]{\scriptsize\color{muted}[100,00\%; 100,00\%]}} & \makecell{100,00\%\\[-1pt]{\scriptsize\color{muted}[100,00\%; 100,00\%]}} & \makecell{100,00\%\\[-1pt]{\scriptsize\color{muted}[100,00\%; 100,00\%]}}\\
Random Forest & BruteForce & \makecell{100,00\%\\[-1pt]{\scriptsize\color{muted}[100,00\%; 100,00\%]}} & \makecell{100,00\%\\[-1pt]{\scriptsize\color{muted}[99,99\%; 100,00\%]}} & \makecell{100,00\%\\[-1pt]{\scriptsize\color{muted}[99,99\%; 100,00\%]}} & \makecell{100,00\%\\[-1pt]{\scriptsize\color{muted}[99,99\%; 100,00\%]}}\\
\rowcolor{rowalt}SVM (RBF) & BruteForce & \makecell{100,00\%\\[-1pt]{\scriptsize\color{muted}[99,99\%; 100,00\%]}} & \makecell{99,99\%\\[-1pt]{\scriptsize\color{muted}[99,98\%; 99,99\%]}} & \makecell{99,98\%\\[-1pt]{\scriptsize\color{muted}[99,95\%; 99,99\%]}} & \makecell{99,99\%\\[-1pt]{\scriptsize\color{muted}[99,98\%; 99,99\%]}}\\
XGBoost & BruteForce & \makecell{100,00\%\\[-1pt]{\scriptsize\color{muted}[100,00\%; 100,00\%]}} & \makecell{100,00\%\\[-1pt]{\scriptsize\color{muted}[99,99\%; 100,00\%]}} & \makecell{100,00\%\\[-1pt]{\scriptsize\color{muted}[99,99\%; 100,00\%]}} & \makecell{100,00\%\\[-1pt]{\scriptsize\color{muted}[99,99\%; 100,00\%]}}\\
\bottomrule
\end{tabular}}%
\par\smallskip{\footnotesize\color{muted}В квадратни скоби: 95\% перцентилен bootstrap интервал по формула (\ref{eq:bootstrap}), 1000 повторни извадки.}

\vspace{2pt}
\interprettab{Всяка клетка съдържа стойността на метриката и интервала, в който тя би попаднала в 95\% от повторенията на теста с нова извадка със същия размер (\cref{eq:bootstrap}). Интервалите са много тесни: полуширината за F1 е най-много $\pm\val{SVM-Bot-F1-hw}$ пункта (SVM при Bot).}
\end{table}

\begin{table}[H]
\centering\small
\caption{Сравнение на алгоритмите по двойки: тест на McNemar и сдвоени bootstrap разлики}\label{tab:mcnemar}
\fitbox{%
\begin{tabular}{@{}llrrrrrrrr@{}}
\toprule
\headrow \hd{Атака} & \hd{Двойка\\A / B} & \hd{И двата\\верни} & \hd{Само A\\верен ($b$)} & \hd{Само B\\верен ($c$)} & \hd{И двата\\грешни} & \hd{$\chi^2$} & \hd{$p$} & \hd{$\Delta$F1} & \hd{$\Delta$Sensitivity}\\
\midrule
DDoS & RF / SVM & 59\,970 & 30 & 0 & 0 & 28,03 & $1{,}2\cdot10^{-7}$ & \makecell{0,050\\[-1pt]{\scriptsize\color{muted}[0,033; 0,068]}} & \makecell{0,040\\[-1pt]{\scriptsize\color{muted}[0,020; 0,063]}}\\
\rowcolor{rowalt}DDoS & RF / XGB & 60\,000 & 0 & 0 & 0 & 0,00 & $1{,}0000$ & \makecell{0,000\\[-1pt]{\scriptsize\color{muted}[0,000; 0,000]}} & \makecell{0,000\\[-1pt]{\scriptsize\color{muted}[0,000; 0,000]}}\\
DDoS & SVM / XGB & 59\,970 & 0 & 30 & 0 & 28,03 & $1{,}2\cdot10^{-7}$ & \makecell{-0,050\\[-1pt]{\scriptsize\color{muted}[-0,068; -0,033]}} & \makecell{-0,040\\[-1pt]{\scriptsize\color{muted}[-0,063; -0,020]}}\\
\rowcolor{rowalt}BruteForce & RF / SVM & 56\,449 & 7 & 0 & 1 & 0,00 & $0{,}0156$ & \makecell{0,012\\[-1pt]{\scriptsize\color{muted}[0,005; 0,023]}} & \makecell{0,021\\[-1pt]{\scriptsize\color{muted}[0,007; 0,039]}}\\
BruteForce & RF / XGB & 56\,456 & 0 & 0 & 1 & 0,00 & $1{,}0000$ & \makecell{0,000\\[-1pt]{\scriptsize\color{muted}[0,000; 0,000]}} & \makecell{0,000\\[-1pt]{\scriptsize\color{muted}[0,000; 0,000]}}\\
\rowcolor{rowalt}BruteForce & SVM / XGB & 56\,449 & 0 & 7 & 1 & 0,00 & $0{,}0156$ & \makecell{-0,012\\[-1pt]{\scriptsize\color{muted}[-0,023; -0,005]}} & \makecell{-0,021\\[-1pt]{\scriptsize\color{muted}[-0,039; -0,007]}}\\
PortScan & RF / SVM & 1\,003 & 5 & 1 & 1 & 1,00 & $0{,}2188$ & \makecell{0,397\\[-1pt]{\scriptsize\color{muted}[0,000; 0,930]}} & \makecell{0,594\\[-1pt]{\scriptsize\color{muted}[0,000; 1,378]}}\\
\rowcolor{rowalt}PortScan & RF / XGB & 1\,007 & 1 & 0 & 2 & 0,00 & $1{,}0000$ & \makecell{0,099\\[-1pt]{\scriptsize\color{muted}[0,000; 0,304]}} & \makecell{0,000\\[-1pt]{\scriptsize\color{muted}[0,000; 0,000]}}\\
PortScan & SVM / XGB & 1\,002 & 2 & 5 & 1 & 2,00 & $0{,}4531$ & \makecell{-0,299\\[-1pt]{\scriptsize\color{muted}[-0,894; 0,198]}} & \makecell{-0,594\\[-1pt]{\scriptsize\color{muted}[-1,378; 0,000]}}\\
\bottomrule
\end{tabular}}%
\par\smallskip{\footnotesize\color{muted}$\Delta$ = A $-$ B в процентни пунктове, в скоби 95\% сдвоен bootstrap интервал.}

\vspace{2pt}
\interprettab{Колоните $b$ и $c$ съдържат записите, по които двата модела не са съгласни (\cref{eq:mcnemar}). При Bot SVM греши значимо по-често от RF ($b=\val{mcn-Bot-RF-SVM-b}$, $c=\val{mcn-Bot-RF-SVM-c}$, $p={}$\val{mcn-Bot-RF-SVM-p}) и от XGBoost ($p={}$\val{mcn-Bot-SVM-XGB-p}), а XGBoost е малко по-точен от RF ($b=\val{mcn-Bot-RF-XGB-b}$, $c=\val{mcn-Bot-RF-XGB-c}$, $p={}$\val{mcn-Bot-RF-XGB-p}). При DDoS и BruteForce RF и XGBoost дават напълно еднакви решения, а SVM прави малко повече грешки (\val{mcn-DDoS-RF-SVM-b} при DDoS и \val{mcn-BruteForce-RF-SVM-b} при BruteForce). Тези разлики са статистически значими, но са под 0,05 процентни пункта и нямат практическо значение.}
\end{table}

% ---------------------------------------------------------------------------
\section{Ефект на редукцията на размерността (E1, E5, A4, E4)}\label{sec:r-red}

\begin{table}[H]
\centering\small
\caption{Пълен срещу редуциран вектор на признаците (E1)}\label{tab:e1}
\fitbox{%
\begin{tabular}{@{}llrrrrr@{}}
\toprule
\headrow \hd{Модел} & \hd{Атака} & \hd{F1 пълен} & \hd{F1 редуциран} & \hd{Разлика, пр. п.} & \hd{Обучение, s} & \hd{µs/запис}\\
\midrule
Random Forest & DDoS & 100,000\% & 100,000\% & $+$0,00 & 5,6 / 4,9 & 3,8 / 3,8\\
\rowcolor{rowalt}SVM (RBF) & DDoS & 99,972\% & 99,950\% & $-$0,02 & 2,0 / 1,6 & 56,4 / 11,2\\
XGBoost & DDoS & 100,000\% & 100,000\% & $+$0,00 & 1,3 / 4,1 & 0,5 / 0,6\\
\rowcolor{rowalt}Random Forest & BruteForce & 99,996\% & 99,998\% & $+$0,00 & 1,4 / 3,2 & 4,0 / 4,1\\
SVM (RBF) & BruteForce & 99,977\% & 99,986\% & $+$0,01 & 0,2 / 1,1 & 3,5 / 2,5\\
\rowcolor{rowalt}XGBoost & BruteForce & 100,000\% & 99,998\% & $-$0,00 & 1,3 / 3,0 & 0,6 / 0,6\\
Random Forest & PortScan & 99,703\% & 99,802\% & $+$0,10 & 0,2 / 0,6 & 36,6 / 37,0\\
\rowcolor{rowalt}SVM (RBF) & PortScan & 99,504\% & 99,405\% & $-$0,10 & 0,0 / 0,4 & 5,4 / 4,1\\
XGBoost & PortScan & 99,703\% & 99,703\% & $+$0,00 & 0,4 / 0,7 & 3,0 / 2,4\\
\bottomrule
\end{tabular}}%
\par\smallskip{\footnotesize\color{muted}Времената са във формат „пълен / редуциран“ вектор.}

\vspace{2pt}
\interprettab{Всеки модел е обучен два пъти с едни и същи хиперпараметри: с всички \val{n-features} признака и с филтър и подбор на \val{n-selected} признака. Разликите в F1 са в рамките на $\pm0{,}03$ пункта, т.е. редукцията не губи информация. Тя ускорява класификацията със SVM \val{SVM-DDoS-speedup} пъти при DDoS и \val{SVM-Bot-speedup} пъти при Bot, защото изчисляването на ядрото е пропорционално на размерността.}
\end{table}

\begin{table}[H]
\centering\small
\caption{Подбор на признаци срещу PCA при едни и същи хиперпараметри (E5)}\label{tab:e5}
\fitbox{%
\begin{tabular}{@{}llrrrr@{}}
\toprule
\headrow \hd{Модел} & \hd{Атака} & \hd{PCA компоненти} & \hd{F1 с PCA} & \hd{F1 с подбор} & \hd{Подбор $-$ PCA, пр. п.}\\
\midrule
Random Forest & DDoS & 22 & 100,000\% & 100,000\% & $+$0,00\\
\rowcolor{rowalt}SVM (RBF) & DDoS & 21 & 99,955\% & 99,950\% & $-$0,01\\
XGBoost & DDoS & 22 & 99,998\% & 100,000\% & $+$0,00\\
\rowcolor{rowalt}Random Forest & BruteForce & 22 & 99,993\% & 99,998\% & $+$0,01\\
SVM (RBF) & BruteForce & 20 & 99,972\% & 99,986\% & $+$0,01\\
\rowcolor{rowalt}XGBoost & BruteForce & 22 & 99,991\% & 99,998\% & $+$0,01\\
Random Forest & Bot & 23 & 99,948\% & 99,967\% & $+$0,02\\
\rowcolor{rowalt}SVM (RBF) & Bot & 22 & 99,600\% & 99,597\% & $-$0,00\\
XGBoost & Bot & 23 & 99,942\% & 99,982\% & $+$0,04\\
\bottomrule
\end{tabular}}%

\vspace{2pt}
\interprettab{PCA запазва 95\% от дисперсията с \val{pca-DDoS}--\val{pca-Bot} компонента. Подборът на \val{n-selected} признака постига същото или малко по-високо F1 с по-малко измерения. Освен това запазва оригиналните признаци, които операторът може да тълкува.}
\end{table}

Анализът A4 (\texttt{results/pca\_sweep.csv}) показва, че при BruteForce два признака стигат за 99,9\%, при DDoS пет признака дават \val{ps-select-RF-DDoS-k5}, а при Bot качеството расте до около 20 признака и след това се изравнява. Затова $k=\val{n-selected}$ е разумният компромис.

Времевите признаци на потока (продължителност, интервали между пакетите, периоди на активност и бездействие) зависят от натоварването на мрежата и от машината, на която е уловен трафикът~\cite{arp2022,engelen2021}. Затова в експеримент E4 моделите са обучени повторно без тях. Промяната на F1 е между \val{SVM-Bot-noiat-dF1} (SVM при Bot) и 0,00 пункта, т.е. в рамките на доверителните интервали (\texttt{results/metrics\_without\_IAT.csv}). Моделите разчитат и на други свойства на трафика, като размер на пакетите, TCP флагове и размер на TCP прозореца, и не зависят от времевите признаци.

% ---------------------------------------------------------------------------
\section{Работа при естествено съотношение на класовете (E2, A5)}\label{sec:r-nat}

В реален трафик атаките не са 50\% от записите. В експеримент E2 записаните модели се прилагат върху записи, които не са участвали в обучението, смесени в естественото съотношение на изтеглените файлове: \val{nat-share-DDoS} DDoS, \val{nat-share-BruteForce} BruteForce и \val{nat-share-Bot} Bot (\cref{tab:e2}).

\begin{table}[H]
\centering\small
\caption{Резултати при естествено съотношение на класовете (E2)}\label{tab:e2}
\fitbox{%
\begin{tabular}{@{}llrrrrrrrr@{}}
\toprule
\headrow \hd{Модел} & \hd{Атака} & \hd{Дял атаки} & \hd{Записи} & \hd{TP} & \hd{FP} & \hd{FN} & \hd{Precision} & \hd{Sensitivity} & \hd{Specificity}\\
\midrule
Random Forest & DDoS & 27,7\% & 200\,000 & 55\,495 & 2 & 0 & 100,0\% & 100,0\% & 100,0\%\\
\rowcolor{rowalt}SVM (RBF) & DDoS & 27,7\% & 200\,000 & 55\,478 & 116 & 17 & 99,8\% & 100,0\% & 99,9\%\\
XGBoost & DDoS & 27,7\% & 200\,000 & 55\,495 & 0 & 0 & 100,0\% & 100,0\% & 100,0\%\\
\rowcolor{rowalt}Random Forest & BruteForce & 17,5\% & 160\,925 & 28\,227 & 4 & 1 & 100,0\% & 100,0\% & 100,0\%\\
SVM (RBF) & BruteForce & 17,5\% & 160\,925 & 28\,221 & 12 & 7 & 100,0\% & 100,0\% & 100,0\%\\
\rowcolor{rowalt}XGBoost & BruteForce & 17,5\% & 160\,925 & 28\,227 & 5 & 1 & 100,0\% & 100,0\% & 100,0\%\\
Random Forest & PortScan & 14,0\% & 3\,613 & 504 & 9 & 1 & 98,2\% & 99,8\% & 99,7\%\\
\rowcolor{rowalt}SVM (RBF) & PortScan & 14,0\% & 3\,613 & 501 & 37 & 4 & 93,1\% & 99,2\% & 98,8\%\\
XGBoost & PortScan & 14,0\% & 3\,613 & 504 & 13 & 1 & 97,5\% & 99,8\% & 99,6\%\\
\bottomrule
\end{tabular}}%

\vspace{2pt}
\interprettab{Recall и Specificity остават практически същите като при балансирания тест, защото не зависят от съотношението на класовете. Precision се променя по \cref{eq:baserate}, но остава висока, защото фалшивите аларми са много редки. Най-ниска е при SVM за Bot: \val{SVM-Bot-nat-Precision} (\val{SVM-Bot-nat-FP} фалшиви срещу \val{SVM-Bot-nat-TP} истински аларми). XGBoost при DDoS не вдига нито една фалшива аларма.}
\end{table}

\Cref{tab:ops} показва работните точки за най-трудната атака, \val{hard-attack}, по двата критерия. Праговете по Precision ограничават шума от аларми, а праговете по Recall ограничават пропуснатите атаки, които са най-нежеланата грешка~\cite{borovska2022,axelsson2000}.

\begin{table}[H]
\centering\small
\caption{Работни точки за Bot при естествено съотношение: праг по целева Precision или по целеви Recall}\label{tab:ops}
\fitbox{%
\begin{tabular}{@{}llrrr@{}}
\toprule
\headrow \hd{Модел} & \hd{Критерий за праг} & \hd{Precision} & \hd{Recall} & \hd{Аларми / 10\,000}\\
\midrule
Random Forest & Precision $\geq$ 90\% & 90,0\% & 100,0\% & 1530\\
\rowcolor{rowalt}Random Forest & Precision $\geq$ 95\% & 95,0\% & 100,0\% & 1450\\
Random Forest & Precision $\geq$ 99\% & 99,0\% & 100,0\% & 1392\\
\rowcolor{rowalt}Random Forest & Recall $\geq$ 90\% & 100,0\% & 97,1\% & 1338\\
Random Forest & Recall $\geq$ 95\% & 100,0\% & 97,1\% & 1338\\
\rowcolor{rowalt}SVM (RBF) & Precision $\geq$ 90\% & 90,0\% & 100,0\% & 1531\\
SVM (RBF) & Precision $\geq$ 95\% & 95,0\% & 99,8\% & 1448\\
\rowcolor{rowalt}SVM (RBF) & Precision $\geq$ 99\% & 99,0\% & 99,3\% & 1382\\
SVM (RBF) & Recall $\geq$ 90\% & 100,0\% & 90,0\% & 1241\\
\rowcolor{rowalt}SVM (RBF) & Recall $\geq$ 95\% & 100,0\% & 95,0\% & 1310\\
XGBoost & Precision $\geq$ 90\% & 90,0\% & 100,0\% & 1531\\
\rowcolor{rowalt}XGBoost & Precision $\geq$ 95\% & 95,4\% & 100,0\% & 1444\\
XGBoost & Precision $\geq$ 99\% & 99,0\% & 100,0\% & 1392\\
\rowcolor{rowalt}XGBoost & Recall $\geq$ 90\% & 100,0\% & 90,0\% & 1240\\
XGBoost & Recall $\geq$ 95\% & 100,0\% & 95,0\% & 1309\\
\bottomrule
\end{tabular}}%

\vspace{2pt}
\interprettab{При критерий по Precision се избира прагът с максимален Recall, а при критерий по Recall се избира най-високият праг, който го постига. При RF и XGBoost Precision 99\% се постига с практически пълен Recall (\val{op-RF-Bot-99-recall}). Затова не е нужен компромис между малко фалшиви аларми и много открити атаки.}
\end{table}

Анализът A7 проверява дали модел за една атака открива другите две. Всеки модел открива почти 100\% от своята атака и практически 0\% от другите. Моделите са научили конкретния отпечатък на атаката, а не общо понятие за „необичаен“ трафик, затова практическата IDS се нуждае от отделен модел за всяка категория атаки.

\section{Проверки за приложимост (A10--A14)}\label{sec:r-rob}

Почти перфектните резултати повдигат въпроса дали моделите научават поведението на атаката, или само отпечатъка на конкретния инструмент и случайното разделяне на записите. Петте проверки в този раздел отговарят на този въпрос (\cref{tab:experiments}).

\paragraph{Признаци-отпечатъци (A10).} Началният размер на TCP прозореца (\texttt{Init Fwd Win Byts}, \texttt{Init Bwd Win Byts}) и минималният размер на сегмента (\texttt{Fwd Seg Size Min}) зависят от операционната система и инструмента на изпращача. Без тези три признака точността при случайно разделяне спада най-много с \val{rob-fp-maxdrop} процентни пункта (\cref{fig:rob-split}). Моделите не разчитат само на отпечатъка: атаката се разпознава и по размера на пакетите, броя им и интервалите между тях.

\paragraph{Хронологично разделяне (A11).} Моделът се обучава върху по-ранните 70\% от потоците и се тества върху по-късните 30\%. Границата е в \val{rob-cut-DDoS} при DDoS, \val{rob-cut-BruteForce} при BruteForce и \val{rob-cut-Bot} при Bot. При DDoS и BruteForce всички модели остават над \val{rob-chrono-min-DDoSBF}. При Bot Random Forest и XGBoost спадат до \val{rob-chronological-all-RF-Bot-acc} и \val{rob-chronological-all-XGB-Bot-acc}, защото откриват само \val{rob-chronological-all-XGB-Bot-sens} от по-късните атаки. След 14:42 Bot започва да обменя потоци с различна структура (5 пакета напред и 2 обратно, отговори от около 326 байта), които не се срещат в ранните данни. SVM се справя с тях по-добре: \val{rob-chronological-all-SVM-Bot-acc}, с \val{rob-chronological-all-SVM-Bot-sens} открити атаки.

\begin{figure}[H]
\centering
\includegraphics[width=0.92\textwidth]{24_robustness_splits.png}
\caption{Точност при случайно и хронологично разделяне, с всички признаци и без признаците-отпечатъци (A10, A11, A14)}\label{fig:rob-split}
\vspace{2pt}
\interpret{Всяка клетка е Accuracy, а под нея е делът открити атаки. Премахването на признаците-отпечатъци почти не променя резултата. Хронологичното разделяне е по-трудно, но при DDoS и BruteForce моделите остават над 99,5\%. Единствените червени клетки са Random Forest и XGBoost при Bot, които пропускат около 10\% от по-късните атаки. Комбинираният детектор (последният ред на всяка атака) възстановява откриването до около 100\%.}
\end{figure}

\paragraph{Комбиниран детектор (A14).} Дърветата са най-точни върху познатите модели на атаката, а гладката граница на SVM улавя и новите. Затова е проверен детектор, който обявява атака, ако я открие XGBoost или SVM. При хронологичното разделяне за Bot той постига \val{rob-chronological-all-XGBorSVM-Bot-acc} с \val{rob-chronological-all-XGBorSVM-Bot-sens} открити атаки, при \val{rob-chronological-all-XGBorSVM-Bot-fpr} фалшиви аларми. В основната постановка (същите тестови записи като в \cref{sec:r-op}) цената е повече фалшиви аларми: при Bot Accuracy е \val{hyb-balanced-Bot-acc} (\val{hyb-balanced-Bot-fp} фалшиви и \val{hyb-balanced-Bot-fn} пропуснати), а при естествено съотношение Precision е \val{hyb-natural-Bot-prec}. Комбинацията е подходяща, когато се очакват нови разновидности на познатата атака; когато атаката е добре позната, XGBoost сам е по-прецизен.

\paragraph{Пренос между подвидовете (A12).} Модел, обучен само върху единия инструмент, почти не открива другия (\cref{fig:rob-sub}). HOIC е HTTP наводнение по TCP, а LOIC-UDP е UDP наводнение, затова двете атаки нямат общи признаци. При BruteForce Random Forest открива всички FTP записи, обучен само върху SSH, но без признаците-отпечатъци откриването пада до \val{sub-SSH-FTP-RF-no_fp}. Общият признак е началният TCP прозорец на инструмента Patator, а не поведението на атаката. Изводът за практиката е, че обучаващите данни трябва да съдържат всеки използван инструмент, както е в основния експеримент.

\begin{figure}[H]
\centering
\includegraphics[width=0.85\textwidth]{26_subtype_transfer.png}
\caption{Дял открити атаки от подвид, който не е участвал в обучението (A12)}\label{fig:rob-sub}
\vspace{2pt}
\interpret{Почти всички клетки са близо до 0\%. Изключение е Random Forest при SSH $\to$ FTP с всички признаци, но този пренос се дължи на общия отпечатък на инструмента и изчезва без него. Записите на FTP са само 53 уникални потока, затова редът FTP $\to$ SSH е показан само за пълнота.}
\end{figure}

\paragraph{Нормален трафик от непознат ден (A13).} Нормалният трафик на 21.02.2018 е необичайно еднороден: \val{benign-2102-win65535} от потоците имат начален TCP прозорец 65\,535, а \val{benign-2102-tcp} са по TCP (в другите дни TCP е около \val{benign-other-tcp}, останалото е UDP). Ако DDoS моделът е обучен само с този нормален трафик, той вдига фалшиви аларми за \val{rob-1day-RF-DDoS} (RF) до \val{rob-1day-XGB-DDoS} (XGBoost) от нормалния трафик в другите дни. С нормален трафик от два дни фалшивите аларми в третия, непознат ден падат до \val{rob-2day-RF-DDoS}--\val{rob-2day-SVM-DDoS} (\cref{fig:rob-days}). Затова основният експеримент използва нормален трафик от всички дни. Обратният случай показва границата на подхода: модел за Bot, обучен без нормалния трафик от 21.02, обявява голяма част от него за атака (Random Forest \val{rob-2day-RF-Bot}, XGBoost \val{rob-2day-XGB-Bot}), докато SVM остава на \val{rob-2day-SVM-Bot}.

\begin{figure}[H]
\centering
\includegraphics[width=\textwidth]{25_benign_days.png}
\caption{Фалшиви аларми върху нормален трафик от ден, който не е участвал в обучението (A13)}\label{fig:rob-days}
\vspace{2pt}
\interpret{Светлите стълбове са модели с нормален трафик само от деня на атаката, тъмните са с нормален трафик от два дни. Показан е по-лошият от двата непознати дни; скалата е логаритмична. При DDoS разнообразният нормален трафик намалява фалшивите аларми над сто пъти. При Bot дърветата правят много фалшиви аларми върху необичайния нормален трафик от 21.02, а SVM остава под 0,1\%.}
\end{figure}

Всички допълнителни фигури и таблици (корелационни матрици, разпределения на признаците, чувствителност към хиперпараметрите, стабилност по фолдове, време за обучение, важност на признаците, разновидности на атаките, модели без времеви признаци, пълните резултати от настройката) се генерират в \texttt{results/} на хранилището~\cite{repo}.
